\documentclass[../main.tex]{subfiles}\begin{document}
\section{Metode og process}
\subsection{Process}

En fuld beskrivelse af projektets proces og arbejdsgange findes i Bilag  $\rightarrow$ Procesbeskrivelse.

\subsection{Metoder}
Nedenfor uddybes de primære metoder og værktøjer, der anvendes til udviklingen af systemet.

\subsubsection{Accepttestspecifikation}
Vi anvender Gherkin-syntaksen til projektets accepttests. Gherkin binder forretningskrav og teknisk implementering sammen gennem sin "Givet-Når-Så" struktur, som gør det muligt at definere testbare scenarier i et naturligt sprog. Metoden sikrer, at vores accepttest afgrænses præcist, og at edge cases identificeres og håndteres.

\subsubsection{Arkitektur og design} 
Softwarearkitekturen dokumenteres og visualiseres med en kombination af C4-modellen og N+1-modellen. Det giver en tydelig adskillelse mellem systemets overordnede struktur og den specifikke implementering.

\paragraph{C4:} Skaber et hierarkisk overblik over systemet på forskellige abstraktionsniveauer: Kontekstdiagrammet (interaktion med omverdenen), Containerdiagrammet (de overordnede softwarekomponenter) og Komponentdiagrammet (ansvarsfordelingen internt). Der anvendes diagramnøgler på alle modeller for at sikre en entydig forståelse.

\subparagraph{Kodediagrammer:}
Til at diagrammere koden og gå fra arkitektur til design er der anvendt applikationsmodellering fra N+1 modellen. Fremgangsmåden for applikationsmodellerne har været iterativ og har fulgt følgende metode:
\begin{enumerate}
    \item Konceptuelt sekvensdiagram udarbejdes med indledende tanker vedrørende User Stories implementering. 
    \item Indledende nødvendige klasser identificeres og placeres i tomt klassediagram uden metoder.
    \item Funktionaliteter udvikles ud fra koncept.
    \item Reelt sekvensdiagram opdateres med aktuelle metoder og fundne nødvendige lifelines/klasser.
    \item Fundne nødvendige klasser og metoder placeres i udfyldt klassediagram, der viser den reelle implementation.
\end{enumerate}

Step 3 til 5 gentages indtil produktets funktionalitet gennemfører accepttest og dermed afspejler funktionalitet applikationsmodellen, og det indledende koncept er stadig tydeligt.

\subsubsection{N+1}
Mens C4 beskriver strukturen, bruges N+1 til at dokumentere udvalgte tekniske perspektiver. For at undgå gentagelser fokuserer vi kun på de views, der direkte relaterer sig til softwaredesignet:

\paragraph{Logical View:} 
Applikationsmodeller, der beskriver, hvordan SW Systemet er gået fra Arkitekturens koncept til reel implementering.

\paragraph{Process View:} 
Her beskrives processer og threads, deres ansvar, samarbejde og selve allokeringen af de logiske elementer (lagene, subsystemer osv.).

\paragraph{Deployment View:} 
Den fysiske udrulning af processer og komponenter til noder, og hvor den fysiske netværksforbindelse mellem noderne etableres. Altså beskrives de forskellige kommunikationsprotokoller.

\paragraph{Data View:} 
Her beskrives et overblik over de data, der bruges i systemet. Det kan være data flows, databaser, både nye og allerede eksisterende, tabeller osv.

\paragraph{Security View:} 
Under Security View beskrives systemets sikkerhed. Dette gøres ved hjælp af en gennemgang af de sikkerhedsordninger og punkter i arkitekturen, hvor sikkerhed anvendes, såsom HTTP-autentifikation, databaseautentifikation osv.

\paragraph{Implementation View:} Her repræsenterer modellen den reelle source code samt beskriver de leverancer (Deliverables) og det, der skaber disse. I afsnittet sker desuden en opsummering over den væsentlige organisering af det tidligere nævnte.

\paragraph{Development View:} 
Her opsummeres de informationer, fremtidige udviklere har brug for at vide om opsætningen af udviklingsmiljøet.

\subsubsection{Applikationsmodellering (Fra arkitektur til design)} 
Overgangen fra de overordnede arkitekturtegninger til den reelle kodebase sker iterativt:

- Der udarbejdes et konceptuelt sekvensdiagram, der skitserer forløbet af en User Story.
- De nødvendige klasser identificeres og samles i et tomt klassediagram.
- Funktionaliteten udvikles i koden og verificeres mod accepttestene.
- Efter endt udvikling opdateres sekvens- og klassediagrammerne med de endelige metoder, så de afspejler den faktiske kode.

\subsubsection{Versionsstyringsværktøj og CI/CD}
Projektets kodebase administreres i GitHub. For at sikre kodekvaliteten løbende har vi opsat en CI/CD-pipeline via GitHub Actions, som automatisk håndterer byggeprocesser, linting og softwaretests.\newline

For at beskytte projektets main-branch er der opsat følgende GitHub Rulesets:

- Direkte commits til main er blokeret; al kode skal indlejres via Pull Requests (PR).

- En PR kræver mindst ét godkendt review før den kan merges.

- Alle GitHub Actions skal bestå uden fejl, før koden kan merges.

- Historikken holdes lineær ved kun at tillade "Squash and Rebase" (merge commits er deaktiveret).

- GitHub CodeQL kører statisk analyse på alle PR's for at fange sikkerhedsrisici, dårlig kode og lækkede konfigurationsfiler.

\end{document}
